\documentclass[12pt,a4paper]{report}
\usepackage[utf8]{inputenc}
\usepackage[T1]{fontenc}
\usepackage[french]{babel}
\usepackage{lmodern}
\usepackage{geometry}
\geometry{margin=2.5cm}
\usepackage{setspace}
\onehalfspacing
\usepackage{graphicx}
\usepackage{float}
\usepackage{caption}
\usepackage{subcaption}
\usepackage{hyperref}
\hypersetup{colorlinks=true,linkcolor=blue,urlcolor=blue,citecolor=blue}
\usepackage{fancyhdr}
\pagestyle{fancy}
\fancyhf{}
\lhead{\includegraphics[height=0.9cm]{assets/university_logo.png}}
\rhead{MyClinic — Rapport Universitaire}
\cfoot{\thepage}
\usepackage{titlesec}
\titleformat{\chapter}{\huge\bfseries}{\thechapter}{1em}{}
\usepackage{adjustbox}
\newcommand{\universitylogo}{assets/university_logo.png}
\newcommand{\applogo}{assets/logo.png}
\newcommand{\universityname}{Nom de l'Université}
\newcommand{\departmentname}{Nom du Département}
\newcommand{\reporttitle}{MyClinic — Gestion des Rendez-vous Médicaux}
\newcommand{\authorname}{Votre Nom}
\newcommand{\supervisorname}{Encadrant: Dr. Nom}
\newcommand{\academicyear}{2025/2026}
\newcommand{\reportdate}{\today}
\begin{document}
\begin{titlepage}
\begin{minipage}[t]{0.22\textwidth}
\includegraphics[width=\textwidth]{\universitylogo}
\end{minipage}
\hfill
\begin{minipage}[t]{0.73\textwidth}
\begin{center}
{\Large \universityname}\\[0.2cm]
{\large \departmentname}\\[1.0cm]
{\huge \textbf{\reporttitle}}\\[0.8cm]
\includegraphics[height=2.6cm]{\applogo}\\[0.8cm]
{\large \textbf{Rapport Universitaire}}\\[0.8cm]
\begin{tabular}{rl}
\textbf{Auteur:} & \authorname \\
\textbf{\supervisorname} & \\
\textbf{Année académique:} & \academicyear \\
\textbf{Date:} & \reportdate \\
\end{tabular}
\end{center}
\end{minipage}
\end{titlepage}
\pagenumbering{roman}
\begin{abstract}
Ce rapport présente la conception et le développement de \textbf{MyClinic}, une application mobile pour la gestion des rendez-vous médicaux, réalisée en binôme ou groupe avec répartition des tâches frontend, backend et tests. L’application repose sur \textit{Expo / React Native} et \textit{Firebase} (Authentication, Firestore). Sont détaillés l’analyse des besoins, le cahier des charges, l’architecture, la modélisation UML, l’implémentation, la sécurité, les tests, le déploiement, la gestion de projet (Taiga, Git) et les perspectives.
\end{abstract}
\tableofcontents
\listoffigures
\listoftables
\newpage
\pagenumbering{arabic}
\chapter{Introduction générale}
Contexte de digitalisation du secteur de la santé et objectifs du projet MyClinic. Problèmes des doubles réservations et des pertes d’information, besoin d’une solution mobile fiable et sécurisée.
\chapter{Contexte et Problématique}
Complexité de coordination patient/médecin, contraintes techniques (créneaux, authentification, stockage fiable, UX). Nécessité d’une solution robuste et ergonomique.
\chapter{Analyse des besoins et Cahier des Charges}
\section{Exigences Fonctionnelles}
\begin{itemize}
\item Recherche et filtrage des médecins par spécialité et disponibilité.
\item Prise, modification et annulation de rendez-vous.
\item Consultation des rendez-vous à venir pour patients et médecins.
\item Gestion des disponibilités et planning des médecins.
\item Notifications de confirmation, annulation, rappel.
\end{itemize}
\section{Exigences Non-Fonctionnelles}
\begin{itemize}
\item Performance et réactivité sur mobile.
\item Sécurité et confidentialité des données.
\item Résilience en cas de perte de connexion.
\item Accessibilité et ergonomie cohérentes.
\item Portabilité et maintenabilité du code.
\end{itemize}
\section{Contraintes}
\begin{itemize}
\item Dépendances aux plateformes mobiles (iOS/Android).
\item Connectivité Internet requise pour Firebase.
\item Respect des normes de protection des données.
\item Gestion des rôles (Patient, Médecin, Admin) et permissions associées.
\end{itemize}
\chapter{Architecture et Conception}
\section{Vue d’Ensemble}
Organisation par rôles via Expo Router: espaces Patient, Médecin et Admin. Composants UI réutilisables, thème centralisé, services Firebase pour auth et données.
\section{Frontend}
Expo / React Native: composants \texttt{Card}, \texttt{Badge}, \texttt{Button}; thème \texttt{colors}, \texttt{ThemeProvider}; navigation et état via hooks personnalisés.
\section{Backend}
Firebase Auth pour la gestion des utilisateurs et Firestore pour \texttt{users}, \texttt{doctors}, \texttt{appointments}, \texttt{prescriptions}, \texttt{medicalRecords}, \texttt{notifications}. Transactions pour unicité de créneau et cohérence.
\section{Modules et Services}
Hooks utilitaires: \texttt{useAuth}, \texttt{useRole}, \texttt{useUserDoc}, \texttt{usePaginatedQuery}. Services Firebase pour la logique CRUD et mises à jour atomiques.
\chapter{Modèle de Données (Firestore)}
\section{Collections principales}
\begin{itemize}
\item \texttt{users}: informations personnelles et rôle.
\item \texttt{doctors}: spécialités et planning.
\item \texttt{appointments}: rendez-vous avec identifiant unique \texttt{slotKey}.
\item \texttt{prescriptions}: médicaments et dosage.
\item \texttt{medicalRecords}: historique médical.
\item \texttt{notifications}: alertes pour utilisateurs.
\end{itemize}
\section{Schéma \texttt{appointments}}
\begin{table}[H]
\centering
\caption{Schéma de la collection \texttt{appointments}}
\begin{tabular}{|l|l|l|}
\hline
Champ & Type & Description \\
\hline
id (slotKey) & string & doctorId\_YYYY-MM-DD\_HH:MM \\
doctorId & string & Identifiant du médecin \\
patientId & string & Identifiant du patient \\
doctorName & string & Nom du médecin \\
clinicName & string & Nom de la clinique \\
patientName & string & Nom du patient \\
scheduledAt & Timestamp & Date et heure du rendez-vous \\
status & string & pending|confirmed|cancelled|completed \\
reason & string & Objet du rendez-vous \\
notes & string & Notes additionnelles \\
createdAt & Timestamp & Date de création \\
\hline
\end{tabular}
\end{table}
\chapter{Conception UML}
\section{Diagramme de Classes}
\begin{figure}[H]
\centering
\includegraphics[width=\textwidth]{assets/classes.png}
\caption{Diagramme de classes MyClinic: structure des entités et relations.}
\end{figure}
\section{Diagramme de Cas d’Utilisation}
\begin{figure}[H]
\centering
\adjustbox{max width=\textwidth}{\includegraphics{assets/usecase.png}}
\caption{Cas d’utilisation pour Patient, Médecin, Admin.}
\end{figure}
\section{Diagramme de Séquence}
\begin{figure}[H]
\centering
\includegraphics[width=\textwidth]{assets/sequence.png}
\caption{Séquence: réservation atomique et confirmation médecin.}
\end{figure}
\section{Diagramme d’Activité}
\begin{figure}[H]
\centering
\includegraphics[width=\textwidth]{assets/activity.png}
\caption{Activité: cycle de vie d’un rendez-vous.}
\end{figure}
\chapter{Implémentation détaillée}
\section{Flux de Réservation}
Sélection d’un créneau depuis \texttt{doctors.availability}. Créneaux indisponibles via \texttt{appointments} (statuts \(\neq\) cancelled). Transaction sur \texttt{appointments/\{slotKey\}}; statut initial \texttt{pending}.
\section{Actions Médecin}
Mises à jour: \texttt{confirmed}, \texttt{cancelled}. Rafraîchissement UI et cohérence des données.
\section{Auto-Completion}
Conversion automatique en \texttt{completed} des rendez-vous échus non annulés, au chargement des listes patient et médecin.
\section{Notifications}
Rappels/confirmations (extension possible).
\chapter{Sécurité et Règles}
\section{Authentification}
Gestion des comptes et rôles via Firebase Auth.
\section{Règles Firestore}
Contrôle d’accès par rôle et propriété des documents; unicité par \texttt{slotKey} et logique transactionnelle.
\section{Confidentialité}
Données minimales, pas de secrets exposés côté client.
\chapter{Tests et Validation}
\section{Plan de Tests}
Réservation, collisions, confirmation/annulation, auto-completion, indisponibilités.
\section{Résultats}
Robustesse des transactions, cohérence des statuts, fluidité UX.
\section{Performances}
Temps de réponse, coût des requêtes; index Firestore si nécessaire.
\chapter{Déploiement et Exploitation}
\section{Pré-requis}
Node.js, Expo, Firebase configuré; images dans \texttt{assets/}.
\section{Exécution Locale}
Démarrage mobile; option web si \texttt{react-native-web}; variables Firebase dans \texttt{services/firebase.js}.
\section{Distribution}
Packaging et consignes de publication.
\chapter{Gestion de Projet}
\section{Organisation}
Groupe jusqu’à 4, responsabilités, jalons. Gestion sur Taiga: Kanban, backlog, sprints, tickets.
\section{Suivi}
Avancement hebdomadaire: commits et tâches terminées. Captures Taiga: board, backlog, burndown.
\section{Git}
Lien dépôt GitHub/GitLab; implication de chaque membre visible dans les commits.
\chapter{Résultats et Discussion}
Suppression des doubles réservations, clarté du cycle de vie, efficacité des actions médecin, satisfaction utilisateur; limites et améliorations.
\chapter{Limites et Perspectives}
Scalabilité, notifications push, congés, iCal, analytics, audit.
\chapter{Conclusion}
Synthèse des apports techniques et fonctionnels, conformité aux objectifs, perspectives d’industrialisation.
\chapter*{Références}
\addcontentsline{toc}{chapter}{Références}
\begin{thebibliography}{9}
\bibitem{expo} Expo Documentation
\bibitem{reactnative} React Native Documentation
\bibitem{firebase} Firebase Auth \& Firestore Documentation
\bibitem{uml} UML \& PlantUML Documentation
\bibitem{taiga} Taiga — gestion de projet agile
\end{thebibliography}
\appendix
\chapter{Annexes — Captures d’Écran}
\section{Écrans Patient}
\begin{figure}[H]
\centering
\begin{subfigure}{0.48\textwidth}
\includegraphics[width=\textwidth]{assets/screens/patient_home.png}
\caption{Accueil Patient}
\end{subfigure}
\begin{subfigure}{0.48\textwidth}
\includegraphics[width=\textwidth]{assets/screens/search_doctors.png}
\caption{Recherche de médecins}
\end{subfigure}
\end{figure}
\begin{figure}[H]
\centering
\begin{subfigure}{0.48\textwidth}
\includegraphics[width=\textwidth]{assets/screens/doctor_details.png}
\caption{Détails médecin}
\end{subfigure}
\begin{subfigure}{0.48\textwidth}
\includegraphics[width=\textwidth]{assets/screens/my_appointments.png}
\caption{Mes rendez-vous}
\end{subfigure}
\end{figure}
\section{Écrans Médecin}
\begin{figure}[H]
\centering
\begin{subfigure}{0.48\textwidth}
\includegraphics[width=\textwidth]{assets/screens/doctor_appointments.png}
\caption{Rendez-vous médecin}
\end{subfigure}
\begin{subfigure}{0.48\textwidth}
\includegraphics[width=\textwidth]{assets/screens/availability.png}
\caption{Disponibilités}
\end{subfigure}
\end{figure}
\section{Écrans Admin}
\begin{figure}[H]
\centering
\begin{subfigure}{0.48\textwidth}
\includegraphics[width=\textwidth]{assets/screens/admin_dashboard.png}
\caption{Tableau de bord Admin}
\end{subfigure}
\begin{subfigure}{0.48\textwidth}
\includegraphics[width=\textwidth]{assets/screens/admin_appointments.png}
\caption{Gestion des rendez-vous}
\end{subfigure}
\end{figure}
\end{document}
